% Concerto --- SOSP submission draft (acmart sigconf).
% Compiles standalone in "nonacm" draft mode (Overleaf or a TeX Live install
% with texlive-publishers, which ships acmart + ACM-Reference-Format.bst).
% Requires the shared bibliography pmcsched-refs.bib in the same directory,
% and the pgfgantt and array packages (both ship with TeX Live).
% For a double-blind SOSP submission, replace the \documentclass line with:
%   \documentclass[sigconf,review,anonymous]{acmart}
% and fill in the \acmConference / \setcopyright metadata.

\documentclass[sigconf,nonacm]{acmart}
\settopmatter{printacmref=false}
\setcopyright{none}
\renewcommand\footnotetextcopyrightpermission[1]{}

\usepackage{booktabs}
\usepackage{array}
\usepackage{pgfgantt}

\begin{document}

\title{Concerto: Counter-Driven Co-Optimization of Scheduling and
Cache/Bandwidth Partitioning}

\author{Anonymous Author(s)}
\affiliation{%
  \institution{Affiliation withheld for review}
  \country{}
}

\begin{abstract}
Two schools of thought address shared-resource contention on multicore
servers. The \emph{partitioning} school carves the last-level cache (LLC) and
memory bandwidth into isolated slices; the \emph{scheduling} school argues that
partitioning is unnecessary and reacts to interference by reallocating cores at
microsecond timescales. We argue that these are not competing philosophies but
two actuators on the same underlying quantity---microarchitectural
contention---and that the right system drives \emph{both} from a single
performance-monitoring-counter (PMC) feedback loop. We present the design of
Concerto, a counter-driven controller that jointly chooses thread placement and
Intel RDT cache/bandwidth partitions. Concerto is motivated by a specific gap:
the leading scheduling-only system, Caladan, explicitly cannot eliminate LLC
interference and instead compensates by granting victims extra cores, while
partitioning-only systems react too slowly to phased behavior. By treating
placement and partitioning as a coupled control problem, Concerto closes the LLC
gap without sacrificing the agility of fast core reallocation. We target chiplet
and CXL-attached topologies, where the non-uniform LLC and interconnect break
the static topology heuristics that commodity schedulers rely on, and we
implement Concerto as a \texttt{sched\_ext} BPF scheduler so that the result is a
deployable commodity-kernel artifact rather than a fork.
\end{abstract}

\maketitle

\section{Introduction}
Co-locating latency-critical (LC) and best-effort (BE) workloads is the primary
lever for raising server utilization, but it reintroduces the contention that
isolation was meant to prevent. The systems community has converged on two
answers. \emph{Hardware partitioning}---Cache Allocation Technology (CAT) and
Memory Bandwidth Allocation (MBA) under Intel RDT~\cite{intel-rdt,linux-resctrl},
as used by Heracles~\cite{lo2015heracles} and PARTIES~\cite{chen2019parties}---
slices the memory hierarchy so that a BE task cannot evict an LC task's working
set. \emph{Fast scheduling}---exemplified by Caladan~\cite{fried2020caladan}---
argues that partitioning is ``neither necessary nor sufficient'' because phased,
bursty workloads change their resource needs faster than partitions can be
re-drawn, and instead reacts by reallocating cores in microseconds.

Both positions are partly right, and that is exactly the problem. Caladan's own
design concedes that it \emph{cannot} address LLC interference directly: when an
antagonist pollutes the shared cache, Caladan cannot un-pollute it, so it
compensates by letting the victim steal additional cores to recover lost compute
capacity~\cite{fried2020caladan}. Partitioning systems can remove the LLC
pollution but pay a reaction-latency penalty and leave cores idle when a slice's
owner is momentarily quiescent. Neither system actuates the knob the other one
holds.

Our thesis is that thread placement and hardware partitioning are
\emph{duals}---two actuators on the single state variable of microarchitectural
contention---and that a controller reading the right PMC vector can move both in
concert, dominating either alone. The historical reason this was never built is
that PMC-in-the-scheduler was considered too costly, too non-portable, and too
hard to make stable; we revisit each of these in light of (i) RDT's
hardware-maintained occupancy and bandwidth counters (CMT/MBM), (ii) the
\texttt{sched\_ext} extensible scheduler class~\cite{linux-schedext}, which lets
us deploy a counter-driven policy without forking the kernel, and (iii) the rise
of chiplet and CXL topologies, where the static cache/NUMA heuristics in today's
schedulers no longer describe the hardware.

\paragraph{Contributions.}
\begin{enumerate}
\item A formulation of scheduling and partitioning as a joint, multi-input
multi-output control problem driven by a shared PMC state estimator, with a
stability argument rather than a heuristic.
\item Concerto, a \texttt{sched\_ext} implementation that co-actuates core
placement and RDT CAT/MBA, with a bounded-overhead counter-multiplexing scheme.
\item An evaluation plan on chiplet (AMD Genoa-class) and CXL-attached hardware
that isolates the contribution of \emph{jointness} from either actuator in
isolation.
\end{enumerate}

\section{Background and Motivation}
\paragraph{Counter signals.}
The contention-aware scheduling literature established that last-level-cache
miss rate is a strong, cheap predictor of co-runner
degradation~\cite{zhuravlev2010contention,blagodurov2010contentionaware}, an
insight that traces back to hardware miss-rate monitors designed expressly to
drive scheduling and partitioning~\cite{suh2002memory} and to cache-fair
scheduling that adjusts timeslices from observed IPC~\cite{fedorova2007isolation}.
Modern RDT exposes Cache Monitoring Technology (per-task LLC occupancy) and
Memory Bandwidth Monitoring directly, so the state Concerto needs---occupancy,
bandwidth, miss rate, IPC, and a stall breakdown---is available without
high-overhead sampling.

\paragraph{Why now: topology heuristics are failing.}
Commodity schedulers place threads using a static description of which cores
share a cache or a memory node~\cite{tam2007clustering}. On monolithic dies this
captured most of the benefit cheaply. On chiplet processors the LLC is split
across multiple cache complexes with non-uniform cross-complex latencies, and on
CXL-attached systems bandwidth contention migrates off-package entirely. The
static topology no longer predicts where contention will actually land, which
reopens the case for a measured, counter-driven signal.

\paragraph{The gap, concretely.}
A planned motivating experiment co-locates a memcached LC service with a
cache-thrashing BE antagonist and sweeps the antagonist's intensity. We expect to
reproduce Caladan's result that core-stealing recovers throughput but not LC
service time once the LLC is polluted, and to show that a single static CAT
partition recovers service time but idles cores during LC quiescence---the two
failure modes Concerto is designed to eliminate simultaneously.

\section{Design}
\paragraph{State estimation.}
Concerto maintains a per-task and per-cache-domain estimate of contention from
the RDT counters plus a small set of core PMCs. Counters are scarce and
multiplexed; we bound estimation overhead by (i) reading hardware-maintained RDT
occupancy/bandwidth, which require no event multiplexing, for the steady signal,
and (ii) time-slicing a small rotating set of core events for the
phase-transition signal, with a configurable cap on per-quantum counter work.

\paragraph{Joint actuation.}
Each control quantum, Concerto solves a small assignment-and-allocation step:
which tasks share a cache domain (placement), and how the CAT/MBA budget is split
across the co-resident tasks (partitioning). Placement handles the coarse,
high-cost decision (which contention domain a task lives in); partitioning
handles the fine, low-cost decision (how the domain's capacity is divided),
including releasing a slice back to BE work the instant the LC owner goes idle so
that partitioning does not strand capacity.

\paragraph{Stability.}
The risk in coupling two actuators is oscillation---placement and partitioning
chasing each other. We treat the loop as a MIMO controller with explicit
time-scale separation: placement runs at the slower load-balancing cadence and
partitioning at the faster quantum, with the partitioner treated as the inner
loop and placement as the outer loop. We will argue stability via this
separation and bound the worst-case actuation rate.

\paragraph{Mechanism.}
Concerto is a \texttt{sched\_ext}~\cite{linux-schedext} BPF scheduler. Placement
decisions are expressed through the scheduler class; partition decisions are
applied through \texttt{resctrl}~\cite{linux-resctrl}. Running in
\texttt{sched\_ext} keeps the artifact loadable on a stock kernel and makes the
overhead and safety story defensible, since the policy is sandboxed and
hot-swappable.

\section{Evaluation Plan}
\paragraph{Hardware.} A monolithic-LLC server (Sapphire-Rapids-class, RDT with
CAT+MBA+CMT+MBM), a chiplet server (Genoa-class, multiple cache complexes), and a
CXL-attached testbed.

\paragraph{Workloads.} LC: memcached, Redis, a search-leaf microservice. BE:
SPEC CPU, a stream/bandwidth antagonist, and a garbage-collected batch job (the
Caladan stressor) to exercise abrupt phase shifts.

\paragraph{Baselines.} Caladan (scheduling-only), PARTIES and a Heracles-style
controller (partitioning-only), stock Linux CFS/EEVDF with static CAT, and
core-pinning. The central comparison is Concerto vs.\ each \emph{single}-actuator
baseline, to attribute wins to jointness rather than to either knob.

\paragraph{Metrics and ablations.} p99/p99.9 LC latency, BE throughput, and
machine utilization at fixed SLO. Ablations: placement-only Concerto,
partitioning-only Concerto, and full Concerto; counter-budget sensitivity; and
behavior across the monolithic/chiplet/CXL spectrum to show where jointness pays
the most.

\section{Related Work}
Concerto sits at the junction of two lines. Contention-aware scheduling uses
counters to separate antagonists across cache
domains~\cite{zhuravlev2010contention,blagodurov2010contentionaware,
fedorova2007isolation,tam2007clustering,suh2002memory}; QoS partitioning carves
the hierarchy to protect LC tasks~\cite{lo2015heracles,chen2019parties}; and
Caladan~\cite{fried2020caladan} argues for fast core reallocation over
partitioning. Concerto's distinction, summarized in Table~\ref{tab:related}, is
that it does not choose a side: it drives placement and partitioning from one PMC
loop and is the first, to our knowledge, to close Caladan's acknowledged LLC gap
by co-actuation rather than compensation.

\begin{table*}[t]
\centering
\footnotesize
\setlength{\tabcolsep}{6pt}
\renewcommand{\arraystretch}{1.3}
\begin{tabular}{@{}>{\raggedright}p{2.6cm} >{\raggedright}p{3.1cm} >{\raggedright}p{3.5cm} l >{\raggedright}p{2.3cm} >{\raggedright\arraybackslash}p{2.2cm}@{}}
\toprule
\textbf{System} & \textbf{Signal} & \textbf{Actuator(s)} & \textbf{Mode} & \textbf{Commodity path?} & \textbf{Target HW} \\
\midrule
Suh et al.\ '02~\cite{suh2002memory}        & Cache miss-rate curves        & Scheduling + partitioning (proposed) & React   & No (custom HW)   & Uniform CMP \\
Fedorova et al.\ '07~\cite{fedorova2007isolation} & IPC (cache-fair)        & Scheduling (timeslice)               & React   & Research kernel  & CMP \\
Zhuravlev et al.\ '10~\cite{zhuravlev2010contention} & LLC miss rate        & Scheduling (placement)               & React   & Research         & Multicore \\
Heracles '15~\cite{lo2015heracles}          & Latency probes, util.\        & Partitioning (CAT, DVFS, power)      & React   & Yes              & Uniform \\
PARTIES '19~\cite{chen2019parties}          & Per-service QoS               & Partitioning (multi-resource)        & React   & Yes              & Uniform \\
Caladan '20~\cite{fried2020caladan}         & Mem.\ bandwidth, req.\ time   & Scheduling (fast core realloc.)      & React   & Kernel module    & Uniform \\
\textbf{Concerto} (this work)               & Full PMC + CMT/MBM            & \textbf{Scheduling + partitioning (joint)} & React & \texttt{sched\_ext} & \textbf{Chiplet / CXL} \\
\textbf{Augur} (companion)                  & PMC \emph{time series}        & Scheduling + page tiering            & \textbf{Predict} & \texttt{sched\_ext} & Tiered / CXL \\
\bottomrule
\end{tabular}
\caption{Where Concerto sits relative to counter-driven scheduling and QoS
partitioning. Prior systems commit to a single actuator and react after
contention appears; Concerto co-actuates scheduling and partitioning from one PMC
loop, and the companion Augur design shifts from reaction to prediction.
``Commodity path?'' indicates whether the technique was shown in a deployable
commodity-kernel path rather than simulation or a research OS.}
\label{tab:related}
\end{table*}

\section{Research Plan and Threats to Validity}
The dominant threat is that Concerto reduces to ``Caladan plus a CAT knob.'' To
refute it, the chiplet/CXL results must show a regime where neither single
actuator suffices and the joint controller is strictly better, and the stability
argument must hold under adversarial phase shifts. A secondary threat is counter
overhead in the scheduler path; the \texttt{sched\_ext} sandbox and the
RDT-first estimation strategy are designed to keep this measurable and bounded.
The 18-month plan is shown in Figure~\ref{fig:gantt}.

\begin{figure*}[t]
\centering
\resizebox{\textwidth}{!}{%
\begin{ganttchart}[
  hgrid, vgrid,
  x unit=0.82cm,
  y unit title=0.6cm,
  y unit chart=0.6cm,
  title height=1,
  bar height=0.6,
  bar/.style={fill=black!22},
  group/.style={fill=black!55},
  milestone/.style={fill=black!78},
  bar label font=\small,
  group label font=\small\bfseries,
  milestone label font=\small\itshape,
  title label font=\small
]{1}{18}
\gantttitle{Month}{18} \\
\gantttitlelist{1,...,18}{1} \\
\ganttgroup{Measurement \& motivation}{1}{3} \\
\ganttbar{Caladan / PARTIES reproduction}{1}{3} \\
\ganttgroup{Mechanism}{2}{9} \\
\ganttbar{PMC + CMT/MBM state estimator}{2}{5} \\
\ganttbar{Joint controller + stability proof}{4}{8} \\
\ganttbar{sched\_ext + resctrl integration}{6}{9} \\
\ganttgroup{Evaluation}{8}{14} \\
\ganttbar{Monolithic-LLC baselines}{8}{10} \\
\ganttbar{Chiplet \& CXL studies}{10}{14} \\
\ganttgroup{Writing}{13}{17} \\
\ganttbar{Draft + ablations}{13}{16} \\
\ganttmilestone{SOSP deadline}{17} \\
\ganttmilestone{Artifact evaluation}{18}
\end{ganttchart}%
}
\caption{Concerto 18-month research plan.}
\label{fig:gantt}
\end{figure*}

% Force the full PMC-driven-scheduling reference set to appear in the
% bibliography even where a given citation is not made in-text. Remove this
% line if you want only the works actually cited above.
\nocite{snavely2000symbiotic,snavely2002priorities,bulpin2005hyperthreading,
tam2007clustering,knauerhase2008osobservations,suh2002memory,fedorova2007isolation,
banikazemi2008pam,zhuravlev2010contention,blagodurov2010contentionaware,
merkel2010resourceconscious,blagodurov2011numa}

\bibliographystyle{ACM-Reference-Format}
\bibliography{pmcsched-refs}

\end{document}
